% GENERATED FILE. Edit canonical lessons, then run npm run generate:books.
% canonical-source-hash: fnv1a64:804a8e859abfb2bd
% canonical-lessons: JA-C149-kiru, JA-C149-nugu, JA-C149-kaburu, JA-C149-youfuku, JA-C149-shitagi

\chapter{Putting Clothes On and Taking Them Off}
\label{ch:ja-putting-clothes-on-and-taking-them-off}
\hlchaptermodality{149}

\begin{chapteropening}
\textbf{By the end of this chapter:} \emph{I can say what I put on and take off.}

Complete the last lesson of chapter 149: I can say what I put on and take off.
\end{chapteropening}

\section[\texorpdfstring{\ja{きる} (kiru)}{きる (kiru)}]{\ja{きる} (kiru) — to wear; to put on (clothes)}

\label{lesson:JA-C149-kiru}



\begin{quote}
Before the new one: say the Japanese for to grow a beard, then the Japanese for make-up.
\end{quote}

\subsection*{You'll want to know: \ja{きる}}
\textbf{\ja{きる}} — \emph{kiru} — \textquotedblleft{}to wear; to put on (clothes)\textquotedblright{}.

\begin{quote}
\textbf{\ja{あさ} \ja{ふくを} \ja{きます。}} — \emph{asa fuku o kimasu} — \textquotedblleft{}I put on my clothes in the morning.\textquotedblright{}
\end{quote}

Polite \textbf{\ja{きます}}, \emph{kimasu}; an ichidan verb. Shoes and trousers go on with \textbf{\ja{はく}}, \emph{haku}.

\begin{cousinweb}
\textbf{\ja{きる}} (\emph{kiru}) is a native Japanese word (wago), inherited rather than borrowed from Chinese; its kanji was attached later as a kun'yomi, and its older origin is not settled (JMdict, KANJIDIC2).
\end{cousinweb}

\begin{grammarlens}[title={What you've built}]
One more word for everyday life at A2.
\end{grammarlens}

\subsection*{Your turn}
\begin{itemize}
\raggedright
  \item \emph{Say it:} \emph{kiru}
  \item \emph{Say it:} \emph{kiru}, once more
  \item \emph{Say it:} \emph{asa fuku o kimasu} — I put on my clothes in the morning.
  \item \emph{Recall:} say the Japanese for to grow a beard, then the Japanese for make-up, then say \emph{kiru} again
\end{itemize}

\begin{tcolorbox}[breakable,colback=teal!4,colframe=teal!35!black,title={Before you move on}]
What does \ja{きる} mean? (\textquotedblleft{}To wear; to put on (clothes)\textquotedblright{}.) Say it once more.
\end{tcolorbox}
\section[\texorpdfstring{\ja{ぬぐ} (nugu)}{ぬぐ (nugu)}]{\ja{ぬぐ} (nugu) — to take off (clothes, shoes)}

\label{lesson:JA-C149-nugu}



\begin{quote}
Before the new one: say the Japanese for make-up, then the Japanese for to wear.
\end{quote}

\subsection*{You'll want to know: \ja{ぬぐ}}
\textbf{\ja{ぬぐ}} — \emph{nugu} — \textquotedblleft{}to take off (clothes, shoes)\textquotedblright{}.

\begin{quote}
\textbf{\ja{うちで} \ja{くつを} \ja{ぬぎます。}} — \emph{uchi de kutsu o nugimasu} — \textquotedblleft{}I take off my shoes at home.\textquotedblright{}
\end{quote}

Polite \textbf{\ja{ぬぎます}}, \emph{nugimasu}; a godan verb.

\begin{cousinweb}
\textbf{\ja{ぬぐ}} (\emph{nugu}) is a native Japanese word (wago), inherited rather than borrowed from Chinese; its kanji was attached later as a kun'yomi, and its older origin is not settled (JMdict, KANJIDIC2).
\end{cousinweb}

\begin{grammarlens}[title={What you've built}]
One more word for everyday life at A2.
\end{grammarlens}

\subsection*{Your turn}
\begin{itemize}
\raggedright
  \item \emph{Say it:} \emph{nugu}
  \item \emph{Say it:} \emph{nugu}, once more
  \item \emph{Say it:} \emph{uchi de kutsu o nugimasu} — I take off my shoes at home.
  \item \emph{Recall:} say the Japanese for make-up, then the Japanese for to wear, then say \emph{nugu} again
\end{itemize}

\begin{tcolorbox}[breakable,colback=teal!4,colframe=teal!35!black,title={Before you move on}]
What does \ja{ぬぐ} mean? (\textquotedblleft{}To take off (clothes, shoes)\textquotedblright{}.) Say it once more.
\end{tcolorbox}
\section[\texorpdfstring{\ja{かぶる} (kaburu)}{かぶる (kaburu)}]{\ja{かぶる} (kaburu) — to wear on the head; to put on (a hat)}

\label{lesson:JA-C149-kaburu}



\begin{quote}
Before the new one: say the Japanese for to wear, then the Japanese for to take off.
\end{quote}

\subsection*{You'll want to know: \ja{かぶる}}
\textbf{\ja{かぶる}} — \emph{kaburu} — \textquotedblleft{}to wear on the head; to put on (a hat)\textquotedblright{}.

\begin{quote}
\textbf{\ja{なつは} \ja{ぼうしを} \ja{かぶります。}} — \emph{natsu wa bōshi o kaburimasu} — \textquotedblleft{}In summer I wear a hat.\textquotedblright{}
\end{quote}

Polite \textbf{\ja{かぶります}}, \emph{kaburimasu}; a godan verb.

\begin{cousinweb}
\textbf{\ja{かぶる}} (\emph{kaburu}) is a native Japanese word (wago), inherited rather than borrowed from Chinese; its kanji was attached later as a kun'yomi, and its older origin is not settled (JMdict, KANJIDIC2).
\end{cousinweb}

\begin{grammarlens}[title={What you've built}]
One more word for everyday life at A2.
\end{grammarlens}

\subsection*{Your turn}
\begin{itemize}
\raggedright
  \item \emph{Say it:} \emph{kaburu}
  \item \emph{Say it:} \emph{kaburu}, once more
  \item \emph{Say it:} \emph{natsu wa bōshi o kaburimasu} — In summer I wear a hat.
  \item \emph{Recall:} say the Japanese for to wear, then the Japanese for to take off, then say \emph{kaburu} again
\end{itemize}

\begin{tcolorbox}[breakable,colback=teal!4,colframe=teal!35!black,title={Before you move on}]
What does \ja{かぶる} mean? (\textquotedblleft{}To wear on the head; to put on (a hat)\textquotedblright{}.) Say it once more.
\end{tcolorbox}
\section[\texorpdfstring{\ja{ようふく} (yōfuku)}{ようふく (yōfuku)}]{\ja{ようふく} (yōfuku) — Western clothes; clothes}

\label{lesson:JA-C149-youfuku}



\begin{quote}
Before the new one: say the Japanese for to take off, then the Japanese for to wear on the head.
\end{quote}

\subsection*{You'll want to know: \ja{ようふく}}
\textbf{\ja{ようふく}} — \emph{yōfuku} — \textquotedblleft{}Western clothes; clothes\textquotedblright{}.

\begin{quote}
\textbf{\ja{ちちは} \ja{まいにち} \ja{ようふくを} \ja{きます。}} — \emph{chichi wa mainichi yōfuku o kimasu} — \textquotedblleft{}My father wears Western clothes every day.\textquotedblright{}
\end{quote}

\begin{cousinweb}
\textbf{\ja{ようふく}} (\emph{yōfuku}) is Sino-Japanese (kango): two Chinese characters read with their on'yomi, readings borrowed from Chinese. They are \emph{yō}, \textquotedblleft{}ocean, Western\textquotedblright{} (Middle Chinese \emph{yang}) and \emph{fuku}, \textquotedblleft{}clothing\textquotedblright{} (\emph{bjuwk}) (KANJIDIC2; Middle Chinese in Baxter's transcription).
\end{cousinweb}

\begin{grammarlens}[title={What you've built}]
One more word for everyday life at A2.
\end{grammarlens}

\subsection*{Your turn}
\begin{itemize}
\raggedright
  \item \emph{Say it:} \emph{yōfuku}
  \item \emph{Say it:} \emph{yōfuku}, once more
  \item \emph{Say it:} \emph{chichi wa mainichi yōfuku o kimasu} — My father wears Western clothes every day.
  \item \emph{Recall:} say the Japanese for to take off, then the Japanese for to wear on the head, then say \emph{yōfuku} again
\end{itemize}

\begin{tcolorbox}[breakable,colback=teal!4,colframe=teal!35!black,title={Before you move on}]
What does \ja{ようふく} mean? (\textquotedblleft{}Western clothes; clothes\textquotedblright{}.) Say it once more.
\end{tcolorbox}
\section[\texorpdfstring{\ja{したぎ} (shitagi)}{したぎ (shitagi)}]{\ja{したぎ} (shitagi) — underwear}

\label{lesson:JA-C149-shitagi}



\begin{quote}
Before the new one: say the Japanese for to wear on the head, then the Japanese for Western clothes.
\end{quote}

\subsection*{You'll want to know: \ja{したぎ}}
\textbf{\ja{したぎ}} — \emph{shitagi} — \textquotedblleft{}underwear\textquotedblright{}.

\begin{quote}
\textbf{\ja{まいにち} \ja{したぎを} \ja{あらいます。}} — \emph{mainichi shitagi o araimasu} — \textquotedblleft{}I wash my underwear every day.\textquotedblright{}
\end{quote}

\begin{cousinweb}
\textbf{\ja{したぎ}} (\emph{shitagi}) is native Japanese, built from \emph{shita}, \textquotedblleft{}under\textquotedblright{}, and \emph{ki}, \textquotedblleft{}wearing\textquotedblright{}, from \emph{kiru}, voiced to \emph{gi}; both parts are inherited rather than borrowed (JMdict, KANJIDIC2).
\end{cousinweb}

\begin{grammarlens}[title={What you've built}]
One more word for everyday life at A2.
\end{grammarlens}

\subsection*{Your turn}
\begin{itemize}
\raggedright
  \item \emph{Say it:} \emph{shitagi}
  \item \emph{Say it:} \emph{shitagi}, once more
  \item \emph{Say it:} \emph{mainichi shitagi o araimasu} — I wash my underwear every day.
  \item \emph{Recall:} say the Japanese for to wear on the head, then the Japanese for Western clothes, then say \emph{shitagi} again
\end{itemize}

\begin{tcolorbox}[breakable,colback=teal!4,colframe=teal!35!black,title={Before you move on}]
What does \ja{したぎ} mean? (\textquotedblleft{}Underwear\textquotedblright{}.) Say it once more.
\end{tcolorbox}
